% CN 9.10 Exercises (p.342--346)

% ======================================================================
\begin{exercise}{1}
    A random word is chosen from a list of all five-letter English words, and
    then one of that word's five letters is chosen.
    \begin{enumerate}[label=(\alph*)]
        \item What is a suitable sample space? What should the probabilities of
              its elements be?
        \item Challenge: use Linux and \texttt{/usr/share/dict/words} to
              determine the probability of each English letter in this
              experiment.
    \end{enumerate}
\end{exercise}

\begin{exercise}{2}
    A fair coin is tossed three times; the outcome is the sequence of results (H
    or T). Specify a suitable sample space, with associated probabilities.
\end{exercise}

\begin{exercise}{3}
    A fair coin is tossed repeatedly until it comes up Heads for the first time;
    the outcome is the sequence of results (e.g.\ TTH). Specify a suitable sample
    space, with associated probabilities.
\end{exercise}

\begin{exercise}{4}
    Two fair dice are thrown and their numbers added; let $T$ be the total.
    \begin{enumerate}[label=(\alph*)]
        \item What is the probability that $T$ is prime?
        \item What is the minimum $d$ such that $\Pr(|T-7|\ge d)\le0.1$?
    \end{enumerate}
\end{exercise}

\begin{exercise}{5}
    Find the probability that a random poker hand contains no diamonds.
\end{exercise}

\begin{exercise}{6}
    A random positive integer $\le100$ is chosen uniformly. What is the
    probability that it is divisible by 3 but not by 9?
\end{exercise}

\begin{exercise}{7}
    (Birthday Paradox) Let $n\in\N$, and suppose $n$ people are chosen at random.
    \begin{enumerate}[label=(\alph*)]
        \item What is the probability that at least two of them share a
              birthday?
        \item What is the minimum $n$ such that it is more likely than not that
              at least two share a birthday?
    \end{enumerate}
\end{exercise}

\begin{exercise}{8}
    Prove that, for any events $A$ and $B$,
    \[
        \Pr(A\cup B) = 1-\Pr(\overline A\cap\overline B).
    \]
\end{exercise}

\begin{exercise}{9}
    Prove by induction on $n$ that, for all $n\in\N$ and any mutually disjoint
    sets $A_1,\dots,A_n$,
    \[
        \Pr(A_1\sqcup A_2\sqcup\cdots\sqcup A_n) = \sum_{k=1}^n\Pr(A_k).
    \]
\end{exercise}

\begin{exercise}{10}
    Suppose events $A,B,C$ are independent. Must $A$ also be independent of
    $B\cap C$? Justify with a proof or a counterexample.
\end{exercise}

% ======================================================================
\begin{exercise}{11}
    Recall the network reliability problems in \S9.7. All links behave
    independently with survival probability $p$. Consider the network
    \begin{center}
    \begin{tikzpicture}[scale=0.9,every node/.style={circle,fill,inner sep=1.5pt}]
        \node[label=left:$S$] (S) at (0,0) {};
        \node[label=above:$M$] (M) at (2,1) {};
        \node[label=below:$N$] (N) at (2,-1) {};
        \node[label=right:$T$] (T) at (4,0) {};
        \draw (S)--(M)--(T)--(N)--(S);
        \begin{scope}[xshift=6.5cm]
            \node[label=left:$S$] (S2) at (0,0) {};
            \node[label=above:$M$] (M2) at (2,1) {};
            \node[label=below:$N$] (N2) at (2,-1) {};
            \node[label=right:$T$] (T2) at (4,0) {};
            \draw (S2)--(M2)--(T2)--(N2)--(S2);
            \draw (M2)--(N2);
        \end{scope}
    \end{tikzpicture}
    \end{center}
    \begin{enumerate}[label=(\alph*)]
        \item (Left network.) Give an expression in $p$ for the probability that
              there is a path of surviving links from $S$ to $T$.
        \item Compare your answer to (a) with Exercise 4.8(b). Discuss the
              relationship.
        \item (Right network: add a link between $M$ and $N$.) Give an expression
              in $p$ for the probability that there is a path of surviving links
              from $S$ to $T$.
        \item Compare your answer to (c) with Exercise 4.8(c). Discuss the
              relationship.
    \end{enumerate}
\end{exercise}

\begin{exercise}{12}
    Determine the probability that a random permutation on a set of two elements
    is fixed-point-free. Then do the same for three elements, and then five.
\end{exercise}

\begin{exercise}{13}
    What is the probability that a Scrabble letter is a vowel, given that it is
    in the first half of the alphabet?
\end{exercise}

\begin{exercise}{14}
    Prove that, for any two events $A$ and $B$, $\Pr(A\mid B)=\Pr(B\mid A)$ if and
    only if the events have the same probability.
\end{exercise}

\begin{exercise}{15}
    Consider again the four-node, five-link network from Exercise 11, with
    $p=\frac12$.
    \begin{enumerate}[label=\arabic*.]
        \item What is the probability that there is a path from $S$ to $T$?
        \item What is the probability that there is a path from $M$ to $N$?
    \end{enumerate}
\end{exercise}

% ======================================================================
\begin{exercise}{16}
    At a cricket match the ball is hit towards a fielder, who could be any of
    Alex, Chris, Kim and Sam (indistinguishable at this distance). Alex catches
    with probability 0.9 (drops with 0.1). Each of Chris, Kim and Sam catches with
    probability 0.4 (drops with 0.6).
    \begin{enumerate}[label=(\alph*)]
        \item What is the probability that the ball is caught?
        \item Given that the ball is caught, what is the probability that the
              fielder was Alex? That it was not Alex?
    \end{enumerate}
\end{exercise}

\begin{exercise}{17}
    A serious crime is committed in a big city; 5 million people cannot be ruled
    out. A degraded DNA fragment from the criminal is found; there is no other
    evidence. A random person has a one-in-a-million chance of matching. A
    database search finds a match with someone who provided DNA in an irrelevant
    context years ago. How likely is it that this person committed the crime? Are
    they guilty beyond reasonable doubt?
\end{exercise}

\begin{exercise}{18}
    In the three-coins example (from p.337 in \S9.9), suppose we observe two
    coin tosses instead of one. For each observation, determine the prior and
    posterior probabilities of the chosen coin being Fair, DoubleHead,
    DoubleTail:
    \begin{enumerate}[label=(\alph*)]
        \item both tosses are Heads;
        \item both tosses are Tails;
        \item the first toss is Heads, the second is Tails.
    \end{enumerate}
\end{exercise}
